% USENIX Security draft. Style: usenix2019_v3 (replace with the year's official
% file when the CfP is out). Body limit at USENIX Security: 13 pages excluding
% references and appendices. See PLAN.md for the evidence map.
\documentclass[letterpaper,twocolumn,10pt]{article}
\usepackage[hidelinks]{hyperref}
\usepackage{usenix2019_v3}
\usepackage{amsmath,amssymb,amsthm}
\usepackage{booktabs}
\usepackage{graphicx}
\usepackage{xspace}
\usepackage[capitalise,noabbrev]{cleveref}
\newcommand{\sys}{TGSA\xspace}

\begin{document}
\date{}
\title{\Large \bf What Log Reduction Costs an Investigation:\\
Lossless Provenance Storage at Collection Rate}
\author{{\rm Anonymous submission}}
\maketitle

\begin{abstract}
% The claim in one paragraph. Draft target: 200 words.
Provenance-based intrusion detection and investigation rest on log
reduction: a host emits far more security-relevant events than any
published system stores, and the reduction applied before storage is
validated only against detection accuracy. We measure what it costs an
investigation. On 58 days of enterprise authentication (1.05 billion
events, 749 labelled red-team actions) and on two whole-system
provenance traces with documented intrusions, causality-preserving
reduction at the ratios these systems operate at deletes up to 44\% of
the attacker's own actions while the field's accuracy metric, recall of
compromised hosts, does not move. We then show that the reduction is
unnecessary: a storage layer that keys provenance on time ingests the
unreduced stream at the rate it is produced, answers forward and
backward causal queries in microseconds to milliseconds, and retains
under a time horizon at the cost of a page release. [Placeholder.]
\end{abstract}

\section{Introduction}
% 1. Provenance graphs are the substrate of attack investigation (BackTracker ... HOLMES, UNICORN).
% 2. Every deployed pipeline reduces the log before storage; reduction is validated by detection accuracy.
% 3. Question: what does reduction cost the investigation, not the detector?
% 4. Findings (E8, E14): attack-edge survival vs node recall; infrastructure reuse erased; granularity.
% 5. The storage pressure that forces reduction is a layout problem (SIGMOD paper, cited as [anonymous]).
% 6. Contributions.

\section{Background and Threat Model}
% provenance graphs; reduction families (CPR, LogGC, NodeMerge, sampling); the investigation workflow; attacker model.

\section{Measuring What Reduction Removes}
% 3.1 metric: attack-edge survival vs node-level recall (E8 Finding 1)
% 3.2 the investigation horizon (E8 Finding 2)
% 3.3 cross-corpus: LANL vs CADETS vs THEIA (E14)
% 3.4 infrastructure reuse (CADETS attack 4)

\section{Lossless Retention at Collection Rate}
% the storage substrate, summarized from the SIGMOD paper: time-keyed addressing,
% forward/backward taint as operators, retention as prefix release; only what the
% security reader needs.

\section{Evaluation}
% E1 ingestion at collection rate (per-host event rates vs engine rate)
% E2 investigation latency on unreduced logs (forward/backward, LANL + TC)
% E3 retention under policy; what a horizon of N days costs
% E4 end-to-end: run a published detector/investigation on reduced vs unreduced input

\section{Discussion: Operating Under Known Contamination}
% the recovery-semantics question from docs/log_reduction_findings.pdf

\section{Related Work}
\section{Conclusion}

\bibliographystyle{plain}
\bibliography{references}
\end{document}
